\section{Energy tails and two-sided truncation}

The off-diagonal estimate and the eigenvector tail of
\cite[Lemmas~6.6--6.7]{Arad2013AreaLaw} have explicit constants.
In this section $P_t=\mathbf1_{(-\infty,t)}(B)$ and $Q_t=I-P_t$.

\begin{lemma}[Spectral separation by imaginary time]
    \label{lem:aklv_spectral_separation}
    \lean{Matrix.IsHermitian.exp_real_smul_eq_cfc,
        Matrix.IsHermitian.norm_spectral_offDiagonal_le_exp_conjugation,
        Matrix.IsHermitian.spectralProjectionGE_mul_one_sub_eq_zero,
        Matrix.IsHermitian.spectral_offDiagonal_add_eq,
        Matrix.IsHermitian.norm_highProjection_mul_exp_neg_le,
        Matrix.IsHermitian.norm_exp_mul_lowProjection_le}
    \leanok
    For Hermitian $B$, any operator $O$, real $t,u$ and $r\ge0$,
    \begin{align}
        \|Q_tOP_u\|&\le e^{-r(t-u)}\|e^{rB}Oe^{-rB}\|,\notag\\
        \|Q_te^{-rB}\|&\le e^{-rt},\qquad
        \|e^{rB}P_u\|\le e^{ru}.
    \end{align}
    If $u\le t$, then $Q_tP_u=0$.
    In that case $Q_t(B+O)P_u=Q_tOP_u$ as well.
\end{lemma}
\begin{proof}\leanok
    Functional calculus gives $e^{rB}=\exp(r\,\cdot)(B)$ and the last
    two estimates by bounding the scalar function on the spectral intervals.
    Insert the inverse exponentials in $Q_tOP_u$ for the first estimate.
    The two spectral indicators have disjoint support when $u\le t$.
\end{proof}

\begin{theorem}[Boundary energy off-diagonal decay]
    \label{thm:aklv_boundary_energy_tail}
    \lean{TNLean.MPS.AreaLaw.AKLV.norm_boundary_spectral_offDiagonal_le,
        TNLean.MPS.AreaLaw.AKLV.norm_two_boundary_spectral_offDiagonal_le,
        TNLean.MPS.AreaLaw.AKLV.norm_exp_conjugation_two_boundary_le_four}
    \leanok
    \uses{def:chain_model}
    Let $0\le h_i\le I$ be nearest-neighbour terms on a chain, let $I_0$
    be any bond set, and put $B=\sum_{i\in I_0}h_i-zI$ for real $z$.
    For any bonds $a,b$ and $u\le t$,
    \begin{align}
        \|Q_th_aP_u\|&\le e^{-(t-u)/8},\notag\\
        \|Q_t(h_a+h_b)P_u\|&\le2e^{-(t-u)/8},\notag\\
        \|e^{B/8}(h_a+h_b)e^{-B/8}\|&\le4.
    \end{align}
    The second inequality is the estimate of
    \cite[Lemma~6.6(2)]{Arad2013AreaLaw}.
\end{theorem}
\begin{proof}\leanok
    \uses{thm:chain_local_imaginary_time,lem:aklv_spectral_separation}
    Center each boundary term as $O_j=h_j-I/2$. Then $\|O_j\|\le1/2$,
    so Theorem~\ref{thm:chain_local_imaginary_time} gives
    $\|e^{B/8}O_je^{-B/8}\|\le1$. Since $Q_tP_u=0$, centering preserves
    the off-diagonal block. Apply Lemma~\ref{lem:aklv_spectral_separation}
    and add the two estimates. Applying the imaginary-time estimate directly
    to $h_a,h_b$, whose norms are at most one, gives the bound $4$.
\end{proof}

\begin{lemma}[Exponential moment of an eigenvector]
    \label{lem:aklv_eigenvector_moment}
    \lean{Matrix.IsHermitian.norm_highProjection_apply_le,
        Matrix.IsHermitian.norm_exp_apply_eigenvector_le,
        Matrix.IsHermitian.norm_highProjection_eigenvector_le_exp}
    \leanok
    Let $B$ be Hermitian and $(B+A)x=ex$ for real $e$.
    Suppose $r\ge0$, $c>0$, $\|e^{rB}Ae^{-rB}\|\le C$, and
    $|e|+C\le c/2$. Then
    \begin{align}
        \|e^{rB}x\|&\le2e^{rc}\|x\|,\notag\\
        \|Q_tx\|&\le2e^{-r(t-c)}\|x\|.
    \end{align}
    Also $\|Q_cy\|\le c^{-1}\|By\|$ for every vector $y$.
\end{lemma}
\begin{proof}\leanok
    \uses{lem:aklv_spectral_separation}
    The function $\mathbf1_{[c,\infty)}(\lambda)/\lambda$ has norm at most
    $c^{-1}$ and its product with $\lambda$ is the high spectral indicator.
    For $y=e^{rB}x$, $By=ey-e^{rB}Ae^{-rB}y$, so
    $\|Q_cy\|\le c^{-1}(|e|+C)\|y\|\le\|y\|/2$.
    Since $\|P_cy\|\le e^{rc}\|x\|$, the triangle inequality bounds
    $\|y\|$ by $2e^{rc}\|x\|$. Multiply by $Q_te^{-rB}$ for the tail.
\end{proof}

\begin{theorem}[Explicit eigenvector energy tail on a chain]
    \label{thm:aklv_ground_energy_tail}
    \lean{TNLean.MPS.AreaLaw.AKLV.norm_groundState_highProjection_le_exp}
    \leanok
    \uses{def:chain_model}
    Under the assumptions of Theorem~\ref{thm:aklv_boundary_energy_tail},
    let $(B+h_a+h_b)x=ex$ with $0\le e\le6$. For every real $t$,
    \begin{align}
        \|Q_tx\|\le2e^{-(t-20)/8}\|x\|.
    \end{align}
    This applies to the ground vector after the regional energy shifts of
    \cite[Section~6.1]{Arad2013AreaLaw} and proves its Lemma~6.7.
\end{theorem}
\begin{proof}\leanok
    \uses{thm:aklv_boundary_energy_tail,lem:aklv_eigenvector_moment}
    Take $r=1/8$, $C=4$ and $c=20$ in
    Lemma~\ref{lem:aklv_eigenvector_moment}, using $|e|+4\le10$.
\end{proof}

\begin{definition}[Separate truncation of the exterior blocks]
    \label{def:aklv_two_sided_truncation}
    \lean{TNLean.MPS.AreaLaw.AKLV.twoSidedEnergyTruncation}
    \leanok
    For exterior operators $L,R$ and a middle operator $M$, put
    \begin{align}
        H^{(t)}=\min(L,tI)+M+\min(R,tI).
    \end{align}
    In \cite[Section~6.1]{Arad2013AreaLaw}, $L=H_L+H_1$,
    $R=H_s+H_R$, and $M=H_2+\cdots+H_{s-1}$.
\end{definition}

\begin{lemma}[Order bounds and agreement below the cutoff]
    \label{lem:aklv_truncation_low_agreement}
    \lean{Matrix.IsHermitian.energyTruncation_mul_lowProjection_eq,
        TNLean.MPS.AreaLaw.AKLV.twoSidedEnergyTruncation_bounds,
        TNLean.MPS.AreaLaw.AKLV.twoSidedEnergyTruncation_mul_lowProjection_eq}
    \leanok
    \uses{def:aklv_two_sided_truncation}
    If $L,M,R\ge0$, $t\ge0$ and $M\le mI$, then
    $0\le H^{(t)}\le L+M+R$ and $H^{(t)}\le(2t+m)I$.
    If $B,L,R$ are Hermitian, $L\le B$, $R\le B$, and both exterior
    blocks commute with $B$, then
    \begin{align}
        H^{(t)}P_t=(L+M+R)P_t.
    \end{align}
    The middle term need not commute with $B$.
\end{lemma}
\begin{proof}\leanok
    Sum the scalar clipping order inequalities. Diagonalize $L,B$
    simultaneously. Their eigenvalues obey $\ell_j\le b_j$, and whenever
    $b_j<t$, $\min(\ell_j,t)=\ell_j$. Thus $\min(L,tI)P_t=LP_t$.
    Repeat for $R$ and add the unchanged middle term.
\end{proof}

% Theorem 6.1 remains open: no completion mark asserts the low-spectrum
% gap or the final phase-aligned squared-distance estimate.
